% ****** Start of file apssamp.tex ******
%
%   This file is part of the APS files in the REVTeX 4.2 distribution.
%   Version 4.2a of REVTeX, December 2014
%
%   Copyright (c) 2014 The American Physical Society.
%
%   See the REVTeX 4 README file for restrictions and more information.
%
% TeX'ing this file requires that you have AMS-LaTeX 2.0 installed
% as well as the rest of the prerequisites for REVTeX 4.2
%
% See the REVTeX 4 README file
% It also requires running BibTeX. The commands are as follows:
%
%  1)  latex apssamp.tex
%  2)  bibtex apssamp
%  3)  latex apssamp.tex
%  4)  latex apssamp.tex
%
\documentclass[%
 reprint,
%superscriptaddress,
%groupedaddress,
%unsortedaddress,
%runinaddress,
%frontmatterverbose, 
%preprint,
%preprintnumbers,
%nofootinbib,
%nobibnotes,
%bibnotes,
 amsmath,amssymb,
 aps,
%pra,
%prb,
%rmp,
%prstab,
%prstper,
%floatfix,
]{revtex4-2}

\usepackage{graphicx}% Include figure files
\graphicspath{{figures/}} % Location of the graphics files
\usepackage{dcolumn}% Align table columns on decimal point
\usepackage{bm}% bold math
%\usepackage{hyperref}% add hypertext capabilities
%\usepackage[mathlines]{lineno}% Enable numbering of text and display math
%\linenumbers\relax % Commence numbering lines

%\usepackage[showframe,%Uncomment any one of the following lines to test 
%%scale=0.7, marginratio={1:1, 2:3}, ignoreall,% default settings
%%text={7in,10in},centering,
%%margin=1.5in,
%%total={6.5in,8.75in}, top=1.2in, left=0.9in, includefoot,
%%height=10in,a5paper,hmargin={3cm,0.8in},
%]{geometry}

\newcommand{\Bzero}{\mathbf{B}_0}
\newcommand{\Uzero}{\mathbf{U}_0}
\newcommand{\vecpot}{\mathbf{A}}
\newcommand{\linop}{\mathcal{L}}
\renewcommand{\Re}{\mathrm{Re}}

\begin{document}

\preprint{APS/123-QED}

\title{Notes on big vs small boxes because slack is getting messy}

%\author{Adrian Fraser}
%\email{adrian_fraser@colorado.edu}
%\affiliation{%
%Department of Astrophysical and Planetary Sciences, University of Colorado, Boulder, CO
%}%

%\author{Alexis K. Kaminski}
%\affiliation{
%Department of Mechanical Engineering, University of California, Berkeley, Berkeley, CA
%}%
%\author{Jeffrey S. Oishi}
%\affiliation{
%Department of Physics and Astronomy, Bates College, Lewiston, ME
%}%

%\date{\today}% It is always \today, today,
             %  but any date may be explicitly specified

%\begin{abstract}
%Shear layers, MHD, optimal perturbations.
%\end{abstract}

%\keywords{Suggested keywords}%Use showkeys class option if keyword
                              %display desired
\maketitle

%\tableofcontents
Below, I'll start with a section on what we've figured out for small boxes, then what we've figured out for big boxes. I'll break each section up into examples of what happens at different Alfv\'{e}n Mach numbers, and I guess at different Reynolds numbers once I get there.

As a reminder, note that the Alfv\'{e}n Mach number is defined as the ratio of the flow speed to the Alfv\'{e}n speed, so the important thing is that it scales like $U_0/B_0$, i.e.,
\begin{equation}
    M_A \equiv \frac{U_0}{v_A} \propto \frac{U_0}{B_0}.
\end{equation}
So a stronger field corresponds to a smaller $M_A$. The system is \textbf{stable} for $M_A \leq M_{A, \mathrm{crit}}$, where
\begin{equation}
M_{A, \mathrm{crit}} \approx 1.05
\end{equation}
(for $\mathrm{Pm} = 1$ anyway, which is all we've looked at). 
Thus, whenever I show results for $M_A = 1$, that means the system is barely stable. 
Note also that, in our non-dimensionalization where we're measuring flow speeds in units of $U_0$, the Alfv\'{e}n speed becomes $1/M_A$, so the frequency of an Alfv\'{e}n wave in the absence of any background flow becomes
\begin{equation}
    \omega_A = \frac{k_x}{M_A}.
\end{equation}
\section{Small boxes}
The following small box runs are scanned over $k_x$ and $k_y$ where $k_y$ is allowed to be zero or positive, $k_x$ is strictly positive, and the minimum nonzero value of $k_x$ and $k_y$ is $0.1$, corresponding to a ``small" box size of $20\pi$.

\subsection{$M_A = 1$}


%\begin{acknowledgments}
%\end{acknowledgments}


%\bibliography{mhd_linopt}% Produces the bibliography via BibTeX.

\end{document}